\par\Needspace{6\baselineskip}\medskip
\SourceLeaf{136}\pdfbookmark[2]{98.}{article-098}
\begin{center}98.\end{center}
The solution of our problem for the ellipse, developed in the preceding articles, could also be extended to the parabola and the hyperbola. We would regard the parabola as an ellipse in which $a$ and $b$ were infinite quantities, $\varphi=90^0$, and finally $E$, $E'$, $g$, and $G=0$. We would likewise regard the hyperbola as an ellipse in which $a$ were negative and $b$, $E$, $E'$, $g$, $G$, and $\varphi$ were imaginary. We prefer, however, to dispense with these assumptions\raisebox{.45ex}{.} and to treat the problem separately for each of the two kinds of conic section. The striking analogy among all three kinds will then become apparent of itself.

For the PARABOLA we retain the symbols $p$, $v$, $v'$, $F$, $f$, $r$, $r'$, and $t$ with the same meanings as above. From the theory of parabolic motion we have
\[
\begin{aligned}
\mathrm{1)}\qquad&\BareRoot{\dfrac p{2r}}=\cos\tfrac12(F-f)\\[.7em]
\mathrm{2)}\qquad&\BareRoot{\dfrac p{2r'}}=\cos\tfrac12(F+f)
\end{aligned}
\]
\[
\begin{aligned}
\frac{2kt}{p^{\frac32}}&=\operatorname{tang}\tfrac12(F+f)-\operatorname{tang}\tfrac12(F-f)\\
&\quad+\tfrac13\operatorname{tang}\tfrac12(F+f)^3-\tfrac13\operatorname{tang}\tfrac12(F-f)^3\\[.5em]
&=\{\operatorname{tang}\tfrac12(F+f)-\operatorname{tang}\tfrac12(F-f)\}\\
&\quad\Bigl\{1+\operatorname{tang}\tfrac12(F+f)\operatorname{tang}\tfrac12(F-f)\\
&\qquad+\tfrac13(\operatorname{tang}\tfrac12(F+f)-\operatorname{tang}\tfrac12(F-f))^2\Bigr\}\\[.5em]
&=\frac{2\sin f\,\mathpunct{.}\BareRoot{rr'}}p
\left\{\frac{2\cos f\,\mathpunct{.}\BareRoot{rr'}}p+\frac{4\sin f^2rr'}{3pp}\right\},
\end{aligned}
\]
\noindent whence
\[
\mathrm{3)}\qquad kt=\frac{2\sin f\cos f\,\mathpunct{.}\,rr'}{\BareRoot{p}}+\frac{4\sin f^3(rr')^{\frac32}}{3p^{\frac32}}\mathpunct{\cdot}
\]
\noindent Further, multiplying equations 1 and 2 gives
\[
\mathrm{4)}\qquad \frac p{\BareRoot{rr'}}=\cos F+\cos f,
\]
\noindent and adding their squares gives
\[
\mathrm{5)}\qquad \frac{p(r+r')}{2rr'}=1+\cos F\cos f.
\]
\noindent Eliminating $\cos F$ from these equations gives
\[
\mathrm{6)}\qquad p=\frac{2rr'\sin f^2}{r+r'-2\cos f\,\mathpunct{.}\BareRoot{rr'}}\mathpunct{\cdot}
\]
\noindent If we therefore adopt equations 9 and $9^*$ of Art. 88 here as well, the former for \SourceLeaf{137}positive $\cos f$ and the latter for negative values, we have
\[
\begin{aligned}
\mathrm{7)}\qquad&p=\frac{\sin f^2\BareRoot{rr'}}{2l\cos f}\\[.6em]
\mathrm{7^*)}\qquad&p=\frac{\sin f^2\BareRoot{rr'}}{-2L\cos f}\mathpunct{\raisebox{.45ex}{,}}
\end{aligned}
\]
\noindent Substituting these values in equation 3 and retaining the symbols $m$ and $M$ with the meanings assigned by equations 11 and $11^*$ of Art. 88 gives
\[
\begin{aligned}
\mathrm{8)}\qquad&m=l^{\frac12}+\tfrac43l^{\frac32}\\[.5em]
\mathrm{8^*)}\qquad&M=-L^{\frac12}+\tfrac43L^{\frac32}.
\end{aligned}
\]
\noindent These equations agree with 12 and $12^*$ of Art. 88 when we set $g=0$ there. It follows that if two heliocentric positions for which a parabola is a solution are treated as if the orbit were elliptical, applying the instructions of Art. 91 must immediately give $x=0$. Conversely, if those instructions give $x=\mathpunct{\raisebox{.65ex}{.}}0$, it is easy to see that the supposed ellipse turns out to be a parabola, since equations 1, 16, 17, 19, and 20 give $b=\infty$, $a=\infty$, and $\varphi=90^0$. It is then very easy to finish determining the elements. For $p$ we may use either equation 7 of this article or equation 18 of Art. 95\footnote{It is also clear from this that $y$ and $Y$ express the same ratios in the parabola as in the ellipse. See Art. 95.}. For $F$, equations 1 and 2 of this article give $\operatorname{tang}\tfrac12 F=\dfrac{\BareRoot{r'}-\BareRoot{r}}{\BareRoot{r'}+\BareRoot{r}}\operatorname{cotang}\tfrac12 f=\sin2\omega\operatorname{cotang}\tfrac12 f$, provided the auxiliary angle has the same meaning as in Art. 89.

We also observe here that if we substitute in equation 3 for $p$ its value from equation 6, we obtain the familiar equation
\[
kt=\tfrac13(r+r'+\cos f\,\mathpunct{.}\BareRoot{rr'})(r+r'-2\cos f\,\mathpunct{.}\BareRoot{rr'})^{\frac12}\BareRoot{2}.
\]
